\documentclass[12pt,a4paper]{report}

\usepackage{relsize}
\usepackage{graphicx}
\usepackage[margin=3cm]{geometry} 
\usepackage{fancyhdr}
\usepackage{titling} 
\usepackage{mathtools,amsthm,amssymb} 
\usepackage[shortlabels,inline]{enumitem}
\usepackage{indentfirst} 
\usepackage{float}
% \usepackage{setspace}
\usepackage{tikz}\usetikzlibrary{arrows.meta,decorations.pathmorphing}
\usepackage{hyperref}\hypersetup{
    colorlinks=true,
    linkcolor=black, 
    filecolor=blue,
    urlcolor=blue,
    citecolor=cyan
}

\theoremstyle{definition}
\newtheorem{theorem}{Theorem}[section]
\newtheorem{lemma}[theorem]{Lemma}
\newtheorem{corollary}[theorem]{Corollary}
\newtheorem{definition}[theorem]{Definition}
\newtheorem{property}[theorem]{Property}
\newtheorem{proposition}[theorem]{Proposition}
\newtheorem{axiom}[theorem]{Axiom}
\newtheorem{example}[theorem]{Example}
\newtheorem{problem}[theorem]{Problem}

\theoremstyle{remark}
\newtheorem*{remark}{Remark}      
\newtheorem*{solution}{Solution}
\newtheorem*{exercise}{Exercise} 
\newtheorem*{claim}{Claim}

\newcommand{\N}{\mathbb{N}}
\newcommand{\Z}{\mathbb{Z}}
\newcommand{\Q}{\mathbb{Q}}
\newcommand{\R}{\mathbb{R}}
\newcommand{\C}{\mathbb{C}}
\newcommand{\F}{\mathbb{F}}
\newcommand{\rank}{\operatorname{rank}}
\newcommand{\nullity}{\operatorname{nullity}}
\newcommand{\la}{\langle}
\newcommand{\ra}{\rangle}

\renewcommand{\headrulewidth}{0.4pt} 
\renewcommand{\footrulewidth}{0.4pt} 
\setlength{\headheight}{15pt}  
\setlength{\parindent}{22pt}   
\pagestyle{fancy}             
\lhead{\rightmark}              
\chead{}                     
\rhead{}          
\lfoot{\theauthor}        
\cfoot{}
\rfoot{\thepage}              
\linespread{1.5}
% \onehalfspacing
\setcounter{tocdepth}{1} % hide subsections in tableofcontents

\title{Discrete Mathematics}
\author{Ray Chien}
\date{\today}

\begin{document}
\maketitle
\tableofcontents

\chapter{Logic}

\section{Propositional Logic}

\begin{definition}[Proposition]
    A \textbf{proposition} is a statement with presice meaning which have a truth value \textbf{true} or \textbf{flase}. 
\end{definition}

\begin{definition}[Operators]
    An \textbf{operator} or \textbf{connective} combines one or more operand express into a larger expresstion.
\begin{table}[h]
    \centering
    \caption{Operator of propositions (Ordered by Precedence)}
    \begin{tabular}{c|c|c|c|c}
        Formal Name & Nickname & Arity & Symbol & Precedence\\
        \hline
        Negation operator & NOT & Unary & $\neg$ & 1\\
        \hline
        Conjunction operator & AND & Binary & $\land$ & 2 \\
        \hline
        Exclusive disjunction operator & XOR & Binary & $\oplus$ & 3 \\
        \hline
        Disjunction operator & OR & Binary & $\lor$ & 4 \\
        \hline
        Implication operator & IMPLIES & Binary & $\rightarrow$ & 5 \\
        \hline
        Biconditional operator & IFF & Binary & $\leftrightarrow$ & 6
    \end{tabular}
\end{table}
\end{definition}



\begin{definition}[Converse, Inverse, and Contrapositive]
    The three statements of an inplication are defined as followed:
    \begin{align*}
        \text{Converse: }& q \rightarrow p\\
        \text{Inverse: }& \neg p \rightarrow \neg q\\
        \text{Contropositive: }& \neg q \rightarrow p. 
    \end{align*}
\end{definition}

The meaning of the contropositive of $\neg q \rightarrow \neg p$ is same as $p \rightarrow q.$ Note that the ``positive'' in 
``contropositive'' actually means ``position'' in English.
\begin{example}
    $p \lor (\neg p \land q) \equiv (p \lor \neg p) \land (p \lor q) \equiv T \land (p \lor q) \equiv p \lor q.$
\end{example}


\begin{definition}[Equivalence]
    Compound proposition p is logically equivalent to compound proposition q, written $p \equiv q$, if the compound proposition $p \leftrightarrow q$ is a tautology.
\end{definition}

\begin{definition}[Tautology]
    A tautology is a compound proposition that is always true.
\end{definition}

For example, $p \lor \neg p$ is a tautology.


\begin{definition}[Contradiction]
    A contradiction is a compound proposition that is always false.
\end{definition}

For example, $p \land \neg p$ is a contradiction.

Some equivence laws are listed as followed:

\begin{theorem}[Some Logical Equivalences]
\mbox{}
\begin{enumerate}
    \item Identity: $p \land T \equiv p \quad p \lor F \equiv F.$
    \item Domination: $p \land F \equiv F \quad p \lor T \equiv T.$
    \item Idempotent: $p \land p \equiv p \lor p \equiv p.$
    \item Double negation: $\neg \neg p \equiv p.$
    \item Associative: $(p \land q) \land r = p \land (q \land r). \quad (p \lor q) \lor r \equiv p \lor (q \lor r).$
    \item Distribution: $p \land (q \lor r) = (p \land q) \lor (p \land r) \quad p \lor (q \land r) = (p \lor q) \land (p \lor r).$
\end{enumerate} 
\end{theorem}

\begin{theorem}[Logical Equivalences Involving Conditional Statements]
    \mbox{}
    \begin{enumerate}
        \item $p \to q \equiv \neg p \lor q.$
        \item $p \to q \equiv \neg q \to \neg p.$
        \item $(p \to q) \land (p \to r) \equiv p \to (q \land r).$
        \item $(p \to r) \land (q \to r) \equiv (p \lor q) \to r.$
        \item $(p \to q) \lor (p \to r) \equiv p \to (q \lor r).$
    \end{enumerate}
\end{theorem}

\begin{theorem}[Logical Equivalences Involving Biconditional Statements]
    \mbox{}
\begin{enumerate}
    \item $p \leftrightarrow q \equiv (p \rightarrow q) \land (q \rightarrow p).$
    \item $p \leftrightarrow q \equiv \neg p \leftrightarrow \neg q.$
    \item $p \leftrightarrow q \equiv (p \land q) \lor (\neg p \land \neg q).$
    \item $\neg(p \leftrightarrow q) \equiv p \leftrightarrow \neg q.$
    \item $p \leftrightarrow q \equiv \neg(p \oplus q).$
\end{enumerate}
\end{theorem}

\begin{theorem}[Logical Equivalences Involving Exclusive Disjunction]
\mbox{}
\begin{enumerate}
    \item $p \oplus q \equiv (p \lor q) \land \neg(p \land q) \equiv (p \land \neg q) \lor (\neg p \land q) \equiv \neg (p \leftrightarrow q).$
\end{enumerate}
\end{theorem}

\section{Predicate Logic}

\begin{definition}[Predicate logic]
    A \textbf{predicate logic} is an extension of propositional logic that permits concisely reasoning about whole classes of entities.
\end{definition}

For example, ``For all $x \in \{1, 2, 3\}, x < 5$'' is a proposition with $x$ as the \textbf{object}, \{1, 2, 3\} as the \textbf{universess of discourse}, and $x < 5$ as the \textbf{predicate} or the \textbf{propositional function} mapping $x \to \{T, F\}.$ Furthermore, ``For all'' is a kind of \textbf{quantifier,} which is defined as followed.

\chapter{Proofs}

\section{Inference and Introduction}

\begin{definition}[Proofs]
    A \textbf{proof} is a correct(well-reasoned, logically valid) and complete(clear, detailed) argument that rigorously and undeniably establishes the truth of a mathematical statement.
\end{definition}

\begin{definition}[Theorems]
    A \textbf{theorem} is a statement that has been proven to be true.
\end{definition}

\begin{definition}[Axioms, Postulates, Hypotheses, Premises]
    An \textbf{axiom, postulate, hypotheses,} or \textbf{premises} is an assumption defining the structure about which we are reasoning.
\end{definition}

\begin{definition}[Rules of Inference]
    A \textbf{rule of inference} is a pattern of logically valid deductions from hypotheses to conclusions.
\end{definition}

\begin{definition}[Lemmas]
    A \textbf{lemma} is a minor theorem used as stepping-stone to proving a major theorem.
\end{definition}

\begin{definition}[Corollaries]
    A \textbf{corollary} is a minor theorem proved as an easy consequence of major theorem.
\end{definition}

\begin{definition}[Conjectures]
    A \textbf{conjecture} is a statement whose truth value has not been proven.
\end{definition}

\begin{definition}[Theories]
    A \textbf{theory} is a set of all theorems that can be proven from a given set of axioms.
\end{definition}

Now we talk about \textbf{inference rules} to help us construct proofs. Inference rules are tautologies allows us to proof the consequence true by antecedents, that is

\[\text{antecedent 1} \land \text{antecedent 2} \land \ldots \rightarrow \text{consequence},\]

\noindent or we write

\begin{center}
    \begin{tabular}{c}
        antecedent 1\\
        antecedent 2\dots\\
        \hline
        $\therefore$ consequence
    \end{tabular} 
\end{center}

Here are several rules:

\noindent 1. Rule of addition
\begin{center}
    \begin{tabular}{c}
        $p$\\
        \hline
        $\therefore p \lor q$
    \end{tabular} 
\end{center}

\noindent 2. Rule of simplification
\begin{center}
    \begin{tabular}{c}
        $p \land q$\\
        \hline
        $\therefore p$
    \end{tabular} 
\end{center}

\noindent 3. Rule of conjunction
\begin{center}
    \begin{tabular}{c}
        $p$\\
        $q$\\
        \hline
        $\therefore p \land q$
    \end{tabular} 
\end{center}

\noindent 4. Rule of modus ponens
\begin{center}
    \begin{tabular}{c}
        $p$\\
        $p \rightarrow q$\\
        \hline
        $\therefore q$
    \end{tabular} 
\end{center}

\noindent 5. Rule of modus tollens (This kind of proof is called \textbf{indirect/ contrapositive proof}, example: $a$ is rational, $b$ is irrational, prove that $a + b$ is irrational.)
\begin{center}
    \begin{tabular}{c}
        $\neg q$\\
        $p \rightarrow q$\\
        \hline
        $\therefore \neg p$
    \end{tabular} 
\end{center}

\noindent 6. Rule of hypothetical syllogism
\begin{center}
    \begin{tabular}{c}
        $p \rightarrow q$\\
        $q \rightarrow r$\\
        \hline
        $\therefore p \rightarrow r$
    \end{tabular} 
\end{center}

\noindent 7. Rule of disjunctive syllogism
\begin{center}
    \begin{tabular}{c}
        $p \lor q$\\
        $\neg p$\\
        \hline
        $\therefore q$
    \end{tabular} 
\end{center}

Another proof sumilar to indirect proof is \textbf{proof by contradiction}, that is, to show that $p$ is true, we assume $\neg p$ is true, and then find a tautology that is $\neg p \rightarrow (q \land \neg q) \equiv F,$ then we know that $p$ must be true. (For example, the proof of ``$\sqrt{2}$ is irrational'')

\section{Methods and Strategy}

A proof of a statement of the form $\exists x P(x)$ is called an \textbf{existence proof}. If the proof demonstrate how actually find or construct a specific element $a$ such that $P(a)$ is true, then it is called a \textbf{constructive proof}. Otherwise, it is called a \textbf{nonconstructive proof}.

\begin{example}[An Constructive Proof]
    Prove that for any integer $n > 0,$ there exists a sequence of $n$ consecutive composite integers.
\end{example}

\begin{proof}
    Consider $(n + 1)! + i + 1$ for $i = 1, 2, 3, \ldots, n$. We know that $i + 1|(n + 1)! + i$.
\end{proof}

\begin{example}[An Non-Constructive Proof]
    Prove tha there are infinitely many prime numbers.
\end{example}

\begin{proof}
    We show that $\forall n \, \exists p \, (p > n) \land (p \, \text{is prime}).$ Consider $n! + 1,$ since $n! + 1 > 1$, so $n! + 1$ is either prime or composite. Using proof by case, if $n! + 1$ is prime, then we find $p = n! +1 > n$. Otherwise, if $n! + 1$ is composite, it must have prime factor $p.$ If $p \leq n$, then $n! + 1$ mod $p = 1$. A contradiction. Therefore, we have $p > n.$ 
\end{proof}

\chapter{Set Theory}

\section{Definition and Relations}

\begin{definition}[Sets]
    A \textbf{set} is a type of structure representing an unordered collection of zero or more distinct objects.
\end{definition}

There are two kinds of notation representing sets. First, we can denot a set $S$ by listing all of its elements in curly braces. For example $\{1, 2, 3\}$. Note that $\{1, 1, 2, 3\}$ and $\{3, 2, 1\}$ have the same meaning because elements of a set is unordered and distinct. Second, a set in the set builder notation represents a set of all elements satisfying a proposition over a domain. For example, $\{x | x > 0\}$ over $\mathbb{R}.$

\begin{definition}[Equality of Sets]
    Two sets are declared to be \textbf{equal} if and only if they contain excactly the same elements.
\end{definition}

Some symbles represents spcial infite sets, including
\begin{enumerate}
    \item Natual numbers: $\mathbb{N} = \{0, 1, 2, \ldots\}.$
    \item Integers: $\mathbb Z.$
    \item Real numbers: $\mathbb R.$
    \item Empty set, or null: $\varnothing$
\end{enumerate}

\begin{definition}[The Empty Set, or Null]
    $\varnothing$ represents the \textbf{empty set}, or \textbf{null}, which is the unique set that contains no elements.
\end{definition}

Note that $\varnothing = \{\} = \{x | F\}.$ Now lets look about relations. A basic set relations is \textbf{member of}. We write $x \in S$ if $x$ is an element or member of set $S$. Now we can define set equality in terms of $\in$ relation:

\[\forall S, T: (S = T) \leftrightarrow (\forall x: x \in S \leftrightarrow x \in T).\]

Also, we have 

\[\neg \exists x: x \in \varnothing.\]

\begin{definition}[Subsets, Supersets]
    Let $S, T$ be two sets. If $\forall x: x \in S \to x \in T$, we say that $S$ is a \textbf{subset} of $T$, denoted $S \subseteq T$, and $T$ is a \textbf{superset} of $S$, denoted $T \supseteq S.$ Furthermore, $\not\subseteq$ and $\not\supseteq$ denotes the negative statements.
\end{definition}

\begin{theorem}
    For any set $S$, $S \subseteq S.$
\end{theorem}

\begin{proof}
    $\forall x: x \in S \to x \in S$ is true.
\end{proof}

\begin{theorem}
    For any set $S$, $\varnothing \subseteq S.$ 
\end{theorem}

\begin{proof}
    Since $\forall x: x \in \varnothing$ is false, therefore, $\forall x: x \in \varnothing \to x \in S$ is true.
\end{proof}

\begin{definition}[Strict/Proper Subets and Supersets]
    Let $S, T$ be two sets. If $S \subseteq T \land S \neq T$, we say that $S$ is a \textbf{proper/strict subset} of $T$, denoted $S \subset T$, and $T$ is a \textbf{proper/strict superset} of $S$, denoted $T \supset S.$ Furthermore, $\not\subset$ and $\not\supset$ denotes the negative statements.
\end{definition}

\begin{definition}[Cardinality]
    The \textbf{cardinality} of a set $S$, denoted $|S|$, is a measure of how many different elements $S$ has.
\end{definition}

For example, $|\varnothing| = 0, |\{1, 2\}| = 2.$ If $|S| \in \mathbb{N},$ we say that $S$ is \textbf{finite}, otherwise, we say that $S$ is \textbf{infinite}.

\begin{definition}[Power Sets]
    Let $S$ be a set, the \textbf{power set} of $S$
    \[P(S) \coloneqq \{x | x \subset S\}.\]
\end{definition}
\end{document}
